\centerline{\bf \Huge Счёт уголков.}

\begin{flushright}
        \scriptsize{— Не понимаю, что мы делаем не так?\\
                    — А я ничего не понимаю.\\
                    Патрик Стар и Спанч Боб.}
\end{flushright}

\problem{1}
На гипотенузе $BC$ прямоугольного треугольника $ABC$ выбрана точка $K$ так, что $AB = AK$. Отрезок $AK$ пересекает биссектрису $CL$ в её середине. Найдите острые углы треугольника $ABC$.

\problem{2}
Точка $K$ --- середина гипотенузы $AB$ прямоугольного равнобедренного треугольника $ABC$. Точки $L$ и $M$ выбраны на катетах $BC$ и $AC$ соответственно так, что $BL = CM$. Докажите, что треугольник $LMK$ --- также прямоугольный равнобедренный.

\problem{3}
В выпуклом четырёхугольнике $ABCD$ биссектриса угла $B$ проходит через середину стороны $AD$, а $\angle C = \angle A + \angle D$. Найдите угол $ACD$.

\problem{4}
На боковых сторонах $AB$ и $AC$ равнобедренного треугольника $ABC$ выбраны точки $P$ и $Q$ соответственно так, что $PQ \parallel BC$. На биссектрисах треугольника $ABC$ и $APQ$, исходящих из вершин $B$ и $Q$, выбраны точки $X$ и $Y$ соответственно так, что $XY \parallel BC$. Докажите, что $PX = CY$.

\problem{5}
В трапеции $ABCD$ точка $M$ --- середина основания $AD$. Известно, что $\angle ABD = 90^{circ}$ и $BC = CD$. На отрезке $BD$ выбрана точка $F$ такая, что $\angle BCF = 90^{circ}$. Докажите, что $MF \perp CD$.

\problem{6}
Дан остроугольный треугольник $ABC$. Высота $AA_1$ продолжена за вершину $A$ на отрезок $A_1A_2 = BC$. Высота $CC_1$ продолжена за вершину $C$ на отрезок $C_1C_2 = AB$. Найдите углы треугольника $A_2BC_2$.

\problem{7}
На биссектрисе $AL$ треугольника $ABC$ выбрана точка $D$. Известно, что $\angle BAC = 2\alpha$, $\angle ACD = 3\alpha$, $\angle ACB = 4\alpha$. Докажите, что $BC + CD = AB$.

\problem{8}
На стороне $AC$ треугольника $ABC$ выбрана точка $D$ такая, что $BD = AC$. Медиана $AM$ этого треугольника пересекает отрезок $BD$ в точке $K$. Оказалось, что $DK = DC$. Докажите, что $AM + KM = AB$.
